% Generated draft: requires mathematical and source review.
\documentclass[11pt]{article}
\usepackage[margin=25mm]{geometry}
\usepackage{fontspec}
\setmainfont{Latin Modern Roman}
\usepackage{amsmath,amssymb,mathtools,unicode-math}
\setmathfont{Latin Modern Math}
\usepackage{enumitem,needspace,xurl,hyperref}
\setlength{\parindent}{0pt}
\setlength{\parskip}{5pt}
\setlength{\emergencystretch}{3em}
\newcommand{\N}{\mathbb{N}}
\newcommand{\Z}{\mathbb{Z}}
\newcommand{\Q}{\mathbb{Q}}
\newcommand{\R}{\mathbb{R}}
\newcommand{\C}{\mathbb{C}}
\newcommand{\F}{\mathbb{F}}
\newcommand{\A}{\mathbb{A}}
\newcommand{\PP}{\mathbb{P}}
\newcommand{\E}{\mathbb{E}}
\newcommand{\eps}{\varepsilon}
\newcommand{\dd}{\,\mathrm{d}}
\begin{document}
\section*{UnsolvedMath ID: 30001945}
\textbf{Source problem code:} catalogue entry
\subsection{Definitions and mathematical statement}
Let \(G\) be a group and let \(K\) be any field. For a subgroup \(H\leq G\), define its first Betti number over \(K\) by
\[
b_1(H;K)=\dim_K H_1(H;K),
\]
where group homology uses the trivial \(H\)-action on \(K\). For a finite-index subgroup \(H\), write \([G:H]\) for its index.

The supplied corrected statement asks for a characterization of the groups \(G\) for which
\[
\lim_i\frac{b_1(G_i;K)}{[G:G_i]}
\]
exists for every inverse system of finite-index normal subgroups \(G_i\trianglelefteq G\) satisfying \(\bigcap_iG_i=\{1\}\), and has a value depending only on \(G\), independent of both the system and the field \(K\).

Here an inclusion-based inverse system can be described by a nonempty directed index set \(I\), with \(G_j\subseteq G_i\) whenever \(i\leq j\), and inclusion bonding maps. With this interpretation, the displayed limit is a net limit. If the intended systems are descending sequences, the corresponding requirements are
\[
G=G_0\supseteq G_1\supseteq G_2\supseteq\cdots,
\qquad [G:G_i]<\infty,
\qquad G_i\trianglelefteq G,
\qquad \bigcap_{i\geq0}G_i=\{1\},
\]
and the limit is a sequential limit. The supplied record does not resolve which indexing convention is intended.

Under a finite-real interpretation, the requested common-limit condition would read
\[
\exists L_G\in\R_{\geq0}\quad
\forall K\text{ a field}\quad
\forall (G_i)\text{ admissible},\qquad
\lim_i\frac{\dim_K H_1(G_i;K)}{[G:G_i]}=L_G,
\]
with every numerator required to be finite. This is an explicit possible interpretation, not a finiteness hypothesis supplied by the source: the record assumes neither finite generation nor finite presentation, and does not specify whether infinite Betti numbers or an infinite limiting value are permitted. It also does not specify whether groups admitting no such system should be excluded or satisfy the universal condition vacuously.

The numerator above is the subgroup Betti number \(b_1(G_i;K)\). The historical original-statement field instead contains \(b_1(G/G_i;K)\); the supplied verification metadata explicitly identifies that expression as erroneous and records its correction to the subgroup expression.
\subsection{Short English statement}
Determine which groups have a normalized first-homology limit that exists along every admissible system of finite-index normal subgroups with trivial intersection and is the same for every such system and every field. The supplied correction uses subgroup homology. The indexing convention, treatment of infinite dimensions and limits, and treatment of groups without admissible systems require clarification.
\subsection{Sources}
\begin{itemize}
\item Local registry: \texttt{data/unsolvedmath/problems.json}, permanent UnsolvedMath ID \(30001945\), title ``Chain-Independent Homology Growth over Arbitrary Fields,'' source problem code \texttt{catalogue entry}. The supplied record belongs to the Oberwolfach Reports Open Problems collection, gives proposed year \(2011\), and names no individual proposer.
\item Supplied report citation: \emph{Stratified Spaces: Joining Analysis, Topology and Geometry} (2012), Oberwolfach Reports, pp.\ 3217--3286, DOI \texttt{10.4171/owr/2011/56}. Supplied source URL and literature-source label ``original OWR report'': \url{https://doi.org/10.4171/owr/2011/56}. The background describes these reports as records of mathematical workshops and their open-problem discussions.
\item Supplied literature-source label ``first L2-Betti approximation'': \url{https://arxiv.org/abs/1206.0474}.
\item Literature assessment dated exactly \texttt{2026-08-21}: the supplied curation labels the problem ``partially solved'' and states that L\"uck approximation settles the finitely presented rational case, while Ershov--L\"uck establish important characteristic-\(p\) inequalities; it describes the general independence question across groups, systems, and fields as remaining open. This is supplied dated curation derived from the OWR material, not an independently verified literature assessment. The source itself requests independent verification of time-sensitive status and cited claims.
\item Supplied statement metadata: \texttt{corrected\_verified}. Its verification note says that the quotient-group numerator was repaired to the subgroup Betti number shown in the official report. This records the supplied provenance claim; no independent source inspection is asserted here.
\end{itemize}
\end{document}
